\subsubsection{LIS O(n log n)}
\begin{lstlisting}[style=mystyle, language=C++]
// TC: O(N log N)
// usa: longitud de la LIS estrictamente creciente
#include <bits/stdc++.h>
using namespace std;

int lisLength(const vector<int>& a){
	vector<int> dp;
	for(int x : a){
		auto it = lower_bound(dp.begin(), dp.end(), x);
		if(it == dp.end()) dp.push_back(x);
		else *it = x;
	}
	return (int)dp.size();
}

// TC: O(N log N),   memoria O(N)
// usa: reconstruye una LIS (la que el algoritmo encontro); una valida cualquiera
//   no decreciente (LNDS): cambiar lower_bound -> upper_bound
//   reconstruye via `parent[i]` (forward) + `dpIdx` (posicion final de la cola)
vector<int> lisSequence(const vector<int>& a){
	int n = (int)a.size();
	if(n == 0) return {};
	vector<int> dp;              // dp[k] = menor cola posible de una IS de long k+1
	vector<int> dpIdx;           // dpIdx[k] = indice en a de esa cola
	vector<int> parent(n, -1);   // parent[i] = indice previo en la LIS que termina en i
	for(int i = 0; i < n; i++){
		int pos = (int)(lower_bound(dp.begin(), dp.end(), a[i]) - dp.begin());
		if(pos == (int)dp.size()){
			dp.push_back(a[i]);
			dpIdx.push_back(i);
		}else{
			dp[pos] = a[i];
			dpIdx[pos] = i;
		}
		if(pos > 0) parent[i] = dpIdx[pos - 1];
	}
	// reconstruir desde la cola de la LIS
	vector<int> res;
	for(int i = dpIdx.back(); i != -1; i = parent[i])
		res.push_back(a[i]);
	reverse(res.begin(), res.end());
	return res;
}

// TC: O(N log N),   memoria O(N) (solo un vector extra: n enteros)
// usa: variante compacta. `insertedAt[i]` = pos en dp donde cayo a[i].
//   Reconstruye recorriendo de atras hacia adelante y agarrando
//   el ultimo i con insertedAt[i] == need (need baja de L-1 a 0).
vector<int> lisSequenceCompact(const vector<int>& a){
	int n = (int)a.size();
	if(n == 0) return {};
	vector<int> dp;
	vector<int> insertedAt(n);
	for(int i = 0; i < n; i++){
		int pos = (int)(lower_bound(dp.begin(), dp.end(), a[i]) - dp.begin());
		if(pos == (int)dp.size()) dp.push_back(a[i]);
		else dp[pos] = a[i];
		insertedAt[i] = pos;
	}
	vector<int> res;
	for(int i = n - 1, need = (int)dp.size() - 1; i >= 0 && need >= 0; i--)
		if(insertedAt[i] == need)
			res.push_back(a[i]), need--;
	reverse(res.begin(), res.end());
	return res;
}

// Diferencia entre las 2:
//   lisSequence       usa 2 vectores extra (parent, dpIdx); reconstruye forward.
//   lisSequenceCompact usa 1 vector extra (insertedAt); reconstruye backward.
// Compacta = menos memoria y menos codigo; basta para reconstruir una LIS.
// Si necesitas contar/todas las LIS, o generalizar el "anterior", usa la primera.

int main(){
	vector<int> a = {1, 3, 2, 4, 5, 2};
	cout << lisLength(a) << "\n";            // 4
	auto seq = lisSequence(a);
	for(int x : seq) cout << x << " ";       // p.ej. 1 2 4 5  (alguna LIS valida)
	cout << "\n";
	auto seq2 = lisSequenceCompact(a);
	for(int x : seq2) cout << x << " ";      // p.ej. 1 2 4 5  (alguna LIS valida)
	cout << "\n";
	return 0;
}
\end{lstlisting}
